% chapters/13/checklist.tex
\bagian{Checklist Pencapaian Kompetensi}

Bab ini menuntaskan \SUBCPMK{3.2} dan menutup \CPMK{3}. Centang setiap butir
berikut hanya apabila Anda benar-benar telah melakukannya atau dapat
menjelaskannya.

\begin{checklistbox}
\begin{checklist}
  \item {\bfseries (Indikator a)} Saya dapat menyusun level secara
        \istilah{greyboxing}, dan menjelaskan alasan penyusunannya demikian.
  \item Saya dapat menyusun level dua dimensi dengan \istilah{tilemap} beserta
        \istilah{Tilemap Collider 2D}.
  \item Saya telah menetapkan tinggi lompatan, lebar celah, dan tinggi
        bongkahan sebelum menyusun levelnya.
  \item {\bfseries (Indikator b)} Saya dapat memilih bentuk \istilah{Collider}
        yang tepat bagi setiap benda.
  \item Saya dapat menjelaskan perbedaan \istilah{Collider} biasa dan
        \istilah{trigger}, beserta syarat agar \istilah{trigger} bekerja.
  \item Saya dapat menjelaskan tiga tipe \istilah{Rigidbody} dan menentukan
        tipe yang tepat bagi setiap benda.
  \item Pemain pada level saya dapat berjalan, melompat, mendarat di atas
        bongkahan, dan terhalang oleh dinding.
  \item Saya dapat menjelaskan mengapa masukan dibaca pada \code{Update}
        sedangkan gerak dikerjakan pada \code{FixedUpdate}.
  \item {\bfseries (Indikator c)} Saya dapat menyebutkan empat masalah fisika
        yang paling sering terjadi beserta gejala dan penyebabnya.
  \item Saya dapat memakai panel \istilah{Physics Debug} untuk menampilkan
        bentuk batas benda.
  \item Saya dapat menelusuri masalah fisika menurut urutan yang disarankan,
        yaitu komponen, bentuk batas, pengaturan, baru kode.
  \item Saya telah menimbulkan satu masalah fisika dengan sengaja, mengamati
        gejalanya, lalu memperbaikinya kembali.
  \item Saya dapat menjelaskan mengapa kamera dipindahkan pada
        \code{LateUpdate}, bukan pada \code{Update}.
  \item Catatan praktikum saya memuat masalah fisika yang saya temui beserta
        perbaikannya.
\end{checklist}
\end{checklistbox}
